\documentclass[11pt]{article}
\usepackage[a4paper,margin=25mm]{geometry}
\usepackage{amsmath,amssymb,amsthm,booktabs,array}
\usepackage[T1]{fontenc}
\usepackage{lmodern}
\usepackage[colorlinks=true,urlcolor=blue,citecolor=blue,linkcolor=blue]{hyperref}
\newtheorem{theorem}{Theorem}
\newtheorem{lemma}[theorem]{Lemma}
\newtheorem{corollary}[theorem]{Corollary}
\newcommand{\mm}{\operatorname{mm}}
\newcommand{\Th}{\Theta}
\title{A twin-free answer to Alfuraidan--Khan--Sauras-Altuzarra,\\
arXiv:2610.26875v1, Section 9, Problem 30 (p.~27):\\
a spectral extension of Knor--Sedlar--\v{S}krekovski (2022)}
\author{Haobo Ma\\\small ChronoAI Pte Ltd, Singapore
\and Wenlin Zhang\\\small National University of Singapore\\\small ChronoAI Pte Ltd, Singapore}
\date{5 November 2026}
\begin{document}
\maketitle
\begin{abstract}
We answer Problem 30 in Section 9 (p.~27) of
Alfuraidan, Khan and Sauras-Altuzarra, \emph{Metric dimension and maximum
eigenvalue multiplicity of graph Laplacians},\linebreak arXiv:2610.26875v1,
and extend the metric-dimension classification of Knor, Sedlar and
\v{S}krekovski (2022) to a classification of both maximum spectral
multiplicities for $\Th_{3,3,r}$, for every $r\geq2$.
For every integer $k\geq1$, the theta graph $G_k=\Th_{3,3,6k+4}$ is connected
and twin-free, and satisfies
$\dim(G_k)=\mm_L(G_k)=2$ and $\mm_Q(G_k)=1$.
Its only repeated Laplacian eigenvalues are $1$ and $3$, each of multiplicity
two. The proof gives an explicit resolving pair and reduces spectral
multiplicities to a second-order recurrence. Distinct graph orders give
pairwise nonisomorphic examples. More generally, both maximum multiplicities
of $\Th_{3,3,r}$ for $r\geq3$ are classified by $r$ modulo six;
the length-two case has a separate quadratic repeated factor.
Attaching an equal pendant path at every vertex of two known base graphs
also gives twin-free infinite families in both separation directions.
\end{abstract}

\section{The family and the question}
For a graph $G$ with adjacency matrix $A$ and diagonal degree matrix $D$, put
$L=D-A$ and $Q=D+A$. Write $\mm_X(G)$ for the largest eigenvalue
multiplicity of $X$, not the multiplicity of its largest eigenvalue.
A set of vertices resolves $G$ if the distances to its members distinguish
all vertices; $\dim(G)$ is the smallest size of such a set.
Two distinct vertices $u,v$ are twins when
$N(u)\setminus\{v\}=N(v)\setminus\{u\}$.

Under the heading ``Open Problems'', Problem 30 of
\cite[Section~9, p.~27]{AKS} asks ``Do there exist infinite families''
of connected twin-free graphs for which exactly one equality holds.
We answer this question affirmatively by constructing connected twin-free
families for which exactly one of $\mm_L=\dim$ and $\mm_Q=\dim$ holds.
The families $\Th_{2,2,r}$ studied in \cite[Theorem~25 and Corollary~26]{AKS}
have twins. Here $\Th_{p,q,r}$ has three internally disjoint paths of
\emph{edge lengths} $p,q,r$ between two terminals.
For the metric dimension of general theta graphs see
Knor, Sedlar and \v{S}krekovski (2022)
\cite[arXiv version, Theorems~8 and~14, pp.~6 and~10]{KSS}; their result already
implies the metric-dimension assertion below. We include a distance
certificate to make the argument self-contained.
Benakli, Froitzheim and Martinez (2025)
\cite[Theorem~4.4, p.~11]{BFM} give a later alternative proof of that
metric classification, using internal-vertex counts rather than edge lengths.
The additional content is the simultaneous multiplicity classification and
its twin-free separation consequence. A related exact formula of
Paone and Paone (2026) \cite[Corollary~8.1, pp.~10--11]{PP}
determines the multiplicity of the particular signless eigenvalue $2$;
Corollary~\ref{cor:two} gives its specialization here.
The six-periodicity at the particular eigenvalues $1,3$ also follows from
Bahmani and Kiani \cite[Corollary~16 and Example~17, pp.~5--6]{BK}:
insert six degree-two vertices in an edge of the long path and take
$n=6$, $k=1,2$ in their reduction formula. This does not exclude a
repeated eigenvalue elsewhere; Section~3.2 supplies that exclusion.
Earlier theta characteristic-polynomial formulas appear in
Liu and Lu \cite[Proposition~2.6 and equations (2.1), (2.3)]{LL},
with internal-vertex parameters $(r-1,2,2)$.

\begin{theorem}\label{thm:main}
Let $k\geq1$, $r=6k+4$, and $G=\Th_{3,3,r}$. Then $G$ is connected and
twin-free, $|V(G)|=r+5$, and
\[
 \dim(G)=\mm_L(G)=2,\qquad \mm_Q(G)=1.
\]
More precisely, the eigenvalues $1$ and $3$ of $L$ each have multiplicity
two and all other eigenvalues of $L$ are simple; every eigenvalue of $Q$
is simple. Consequently these graphs answer Problem 30 in the
Laplacian-equality direction.
\end{theorem}

\section{Geometry and a resolving pair}
Label the short paths $u,a,b,v$ and $u,c,d,v$, and the long path
$y_0,y_1,\ldots,y_r$, where $y_0=u$, $y_r=v$.
Connectivity and the order count follow immediately.
Every cycle uses two complete terminal-to-terminal paths, so its length is
$6$ or $r+3$. Thus there are no triangles or quadrilaterals.
Every vertex has degree at least two. Adjacent twins would have a common
third neighbor and give a triangle; nonadjacent twins would have two
common neighbors and give a quadrilateral. Hence $G$ is twin-free.

The distance pairs to $(a,y_1)$ at the six named short-path vertices are
\[
\begin{array}{c|rrrrrr}
 x&u&v&a&b&c&d\\\hline
 (d(x,a),d(x,y_1))&(1,1)&(2,4)&(0,2)&(1,3)&(2,2)&(3,3).
\end{array}
\]
For $1\leq j\leq r-1$ the pair at $y_j$ is
\begin{equation}\label{eq:distance}
 \bigl(\min(j+1,r-j+2),\ \min(j-1,r-j+4)\bigr).
\end{equation}
Indeed a path from $y_j$ to $a$ must first reach a terminal, giving
the two candidate lengths $j+1,r-j+2$. A path to $y_1$ either stays
on the long path or uses a short terminal-to-terminal path of length
three, giving $j-1,r-j+4$.
The displayed six pairs follow by the same terminal-path comparison.

Write $r=2h$, so $h\geq5$. Formula \eqref{eq:distance} becomes
\[
\begin{cases}
 (j+1,j-1),&1\leq j\leq h,\\
 (h+1,h),&j=h+1,\\
 (h,h+1),&j=h+2,\\
 (2h-j+2,2h-j+4),&h+3\leq j\leq 2h-1.
\end{cases}
\]
The coordinate differences are respectively $2,1,-1,-2$.
Each outer range is injective. The middle pairs have coordinates at
least five, so they are distinct from the six listed pairs.
With first minus second, the six listed differences are
$0,-2,-2,-2,0,0$. Thus no listed pair has difference $2$; the listed pairs of difference
$-2$ are $(2,4),(0,2),(1,3)$, whereas the descending long-path pairs
have first coordinate at least three. There are no collisions.
The pair $(a,y_1)$ therefore resolves $G$.
No single vertex resolves $G$: if distances to a vertex were injective,
a vertex of degree three could have at most two neighbors. Indeed, each
neighbor's distance differs from its own by at most one; equality is
excluded by injectivity, leaving only the values one less and one greater.
This is a contradiction. Hence $\dim(G)=2$.

\section{The spectral recurrence}
Treat both matrices at once as $M_e=D-eA$, with $e=1$ for $L$ and
$e=-1$ for $Q$. For an eigenvalue $\lambda$, put $t=e(2-\lambda)$.
Swapping the two short paths is an involution commuting with $M_e$.
The antisymmetric subspace consists of vectors zero on the terminals
and the long path, with values $(B,C)$ on $(a,b)$ and $(-B,-C)$ on
$(c,d)$. Its matrix is
\[
 \begin{pmatrix}2&-e\\-e&2\end{pmatrix}.
\]
Its eigenvalues are $1,3$, each simple, corresponding to $t=1,-1$.
It remains to bound multiplicities in the symmetric subspace.

For a symmetric eigenvector use values $A,D_0$ at $u,v$, values $B,C$
on each short path, and values $Y_j$ on the long path. The eigenvalue
equations are exactly
\begin{align}
 tB&=A+C,&tC&=B+D_0,\label{eq:short}\\
 Y_0&=A,&Y_r&=D_0,\nonumber\\
 Y_1&=(t+e)A-2B,&Y_{r-1}&=(t+e)D_0-2C,\label{eq:ends}\\
 Y_{j+1}&=tY_j-Y_{j-1}&&(1\leq j\leq r-1).\label{eq:rec}
\end{align}
These equations also suffice for an eigenvector, by checking each
degree-two row and the two terminal rows.

\subsection{The parameters \texorpdfstring{$t=1,-1$}{t = 1, -1}}
At $t=1$ the sequence determined by $(Y_0,Y_1)=(A,s)$ is
\[
 A,s,s-A,-A,-s,A-s,A,s,\ldots;
\]
at $t=-1$ it is
\[
 A,s,-A-s,A,s,-A-s,\ldots.
\]
Since $r=6k+4$, the pairs $(Y_r,Y_{r-1})$ are respectively
$(-s,-A)$ and $(s,A)$. Substitution in \eqref{eq:short}--\eqref{eq:ends}
gives all symmetric solutions in the following table:
\[
\begin{array}{cc|ccccc}
 e&t&A&B&C&D_0&s\\\hline
 1&1&A&A/2&-A/2&-A&A\\
 1&-1&A&-A/2&-A/2&A&A\\
 -1&1&0&0&0&0&0\\
 -1&-1&0&0&0&0&0.
\end{array}
\]
For example, for $t=1$ the short equations give $D_0=-A$ and $B-C=A$;
the long path gives $s=A$. The two terminal equations give
$2B=eA$ and $2C=-eA$, so $A=eA$.
For $t=-1$ they give $D_0=A$, $B+C=-A$, $s=A$ and
$2B=2C=(e-2)A$, hence $(e-1)A=0$.
This proves the table, including that every listed row satisfies all
the equations. The symmetric multiplicity is one for $L$ and zero
for $Q$ at these parameters.

\subsection{All other parameters}
Define polynomials $p_0=1$, $p_1=t$, and
$p_{n+1}=tp_n-p_{n-1}$ for $n\geq1$.
Induction gives both
\begin{equation}\label{eq:cont}
 Y_r=p_{r-1}(t)Y_1-p_{r-2}(t)Y_0,\qquad
 p_n(t)^2+p_{n+1}(t)^2-tp_n(t)p_{n+1}(t)=1.
\end{equation}
For the first identity the initial cases $r=2$ and $r=3$ follow from
\eqref{eq:rec}; for the second the expression at $n+1$ equals that at
$n$, and its value at $n=0$ is one.

Suppose now that $q=t^2-1\ne0$. The short-path equations give
\[
 qB=tA+D_0,\qquad qC=A+tD_0.
\]
Thus $(A,D_0)$ determines $B,C,Y_1$, and then every $Y_j$.
Writing $P=p_{r-1}(t)$ and $R=p_{r-2}(t)$, the left terminal equation
and \eqref{eq:cont} imply
\begin{equation}\label{eq:linear}
 (q+2P)D_0=\bigl(P((t+e)q-2t)-Rq\bigr)A.
\end{equation}
We show that this is a nonzero linear relation on $(A,D_0)$.
Otherwise
\[
 P=-q/2,\qquad R=-\bigl((t+e)q-2t\bigr)/2.
\]
Substitution in $R^2+P^2-tRP=1$ gives
\begin{equation}\label{eq:factor}
 (t-2e)(t^2-1)(t^2+2et-1)=0.
\end{equation}
More explicitly, four times the substituted expression $R^2+P^2-tRP-1$ is
$e(t-2e)(t^2-1)(t^2+2et-1)$, which also verifies the signs.
If $t=2e$, the recurrence gives $P=er$, incompatible with $P=-3/2$.
The remaining possibility is $t^2+2et-1=0$. Reducing the expression
for $R$ modulo this quadratic gives $R=e$.
The quadratic is irreducible over $\mathbb Q$, since its roots
$-e\pm\sqrt2$ are irrational. Hence the rational polynomial
$p_{r-2}(t)-e$, vanishing at either root, vanishes at both.
In particular it vanishes at $T=-e(1+\sqrt2)$, where $|T|>2$.
But for $x>2$ the recurrence implies
$p_{n+1}(x)>p_n(x)\geq1$ for $n\geq0$, beginning with
$p_1(x)>p_0(x)=1$; and $p_n(-x)=(-1)^np_n(x)$.
Since $r-2$ is even and at least two, $p_{r-2}(T)>1$, a contradiction.
Thus \eqref{eq:linear} is nonzero. The injective endpoint map places
the symmetric eigenspace in a line in $\mathbb R^2$, so its dimension
is at most one.

Combining the two invariant subspaces proves all spectral assertions
of Theorem~\ref{thm:main}. In particular each of $1,3$ has one symmetric
and one antisymmetric Laplacian eigenvector, while the signless spectrum
has no repeated eigenvalue. The real spectral theorem ensures that
eigenspace dimensions equal algebraic multiplicities.

\begin{corollary}
The graphs $G_k$, $k\geq1$, are pairwise nonisomorphic connected twin-free
graphs satisfying exactly one of $\mm_L=\dim$ and $\mm_Q=\dim$.
\end{corollary}
\begin{proof}
Their orders $6k+9$ are distinct, and Theorem~\ref{thm:main} gives the
equalities and inequality.
\end{proof}

\subsection{The whole subfamily \texorpdfstring{$\Th_{3,3,r}$ for $r\geq6$}{Theta(3,3,r) for r >= 6}}
\begin{corollary}\label{cor:all-lengths}
For every integer $r\geq6$, $\Th_{3,3,r}$ is connected and twin-free with
metric dimension two. Its repeated eigenvalues are exactly those in the
following table; each listed eigenvalue has multiplicity two, and every
other eigenvalue is simple.
\[
\begin{array}{c|cc|cc}
 r\bmod6&\text{repeated in }L&\text{repeated in }Q&\mm_L&\mm_Q\\\hline
 0&3&1&2&2\\
 1&3&3&2&2\\
 2&\varnothing&\varnothing&1&1\\
 3&1,3&1,3&2&2\\
 4&1,3&\varnothing&2&1\\
 5&\varnothing&\varnothing&1&1
\end{array}
\]
Thus within this subfamily, both equalities hold in residues $0,1,3$;
neither holds in residues $2,5$; and only the Laplacian equality holds
in residue $4$.
\end{corollary}
\begin{proof}
The connectivity and twin argument applies unchanged.
For even $r=2h\geq6$, the distance proof applies with $h\geq3$;
the middle pairs have coordinate differences $1,-1$, absent from
the six short-path pairs.
If $r=2h+1\geq7$, the six short-path pairs and
\eqref{eq:distance} remain valid. The long-path pairs are
$(j+1,j-1)$ for $j\leq h+1$, the single pair $(h+1,h+1)$ at $j=h+2$,
and $(2h+3-j,2h+5-j)$ for $j\geq h+3$.
The two outer differences are $2,-2$, and the middle coordinates are
at least four, so the middle pair differs from $(1,1),(2,2),(3,3)$.
As before there are no collisions, and the degree-three
argument excludes a single landmark.

The nonsingular spectral proof did not need $r\equiv4\pmod6$.
For general $r$, at $t=2e$ its incompatible values are
$P=e^{r-1}r$ and $P=-3/2$. At the quadratic root $T$, the same growth
argument gives $|p_{r-2}(T)|>1$, independently of parity. Thus every
nonsingular symmetric eigenspace still has dimension at most one.

For the singular cases write $s=Y_1$. The short equations reduce to
\[
\begin{array}{c|ccc}
 t&D_0&C&s\\\hline
 1&-A&B-A&(1+e)A-2B\\
 -1&A&-A-B&(e-1)A-2B.
\end{array}
\]
Insert the periodic sequences from Section~3.1 into the two remaining
conditions $Y_r=D_0$ and $Y_{r-1}=(t+e)D_0-2C$.
These are two homogeneous equations in $(A,B)$.
For $t=1$, their coefficient determinants for residues $0,1,2,4,5$
are respectively
\[
8,\quad4(e+1),\quad4e,\quad4(1-e),\quad4(2-e).
\]
In residue $3$ the only constraint is $(2e-1)A=0$.
For $t=-1$, residues $r\equiv0\pmod3$ give the single constraint
$(2e+1)A=0$; residues $1,2$ have determinants
$4(1-e),4(e+2)$ respectively.
Every zero-determinant matrix in this list has a nonzero coefficient
$2$ or $-2$ of $B$, so its rank is one, not zero.
It follows that the symmetric eigenspace dimensions are
\[
\begin{array}{c|cc|cc}
 r\bmod6&L:t=1&L:t=-1&Q:t=1&Q:t=-1\\\hline
 0&0&1&0&1\\
 1&0&1&1&0\\
 2&0&0&0&0\\
 3&1&1&1&1\\
 4&1&1&0&0\\
 5&0&0&0&0.
\end{array}
\]
Adding the one-dimensional antisymmetric eigenspace at each singular
parameter and using $\lambda=2-et$ proves the stated table.
\end{proof}

\begin{corollary}\label{cor:complete}
The graphs $\Th_{3,3,r}$ are connected and twin-free for every $r\geq2$.
Their metric dimension is three for $r=3,5$, and two otherwise.
The spectral table in Corollary~\ref{cor:all-lengths} holds for every
$r\geq3$. For $r=2$, the only repeated Laplacian eigenvalues are
$3-\sqrt2,3+\sqrt2$, both double, and the signless spectrum is simple.
Consequently exactly one equality holds precisely when
\[
 r=2\quad\hbox{or}\quad r\equiv4\pmod6,
\]
and in every such case it is the Laplacian equality.
Both equalities hold precisely for $r\geq6$ with $r\equiv0,1,3\pmod6$.
Neither holds in the other cases. In particular, no graph of this
entire subfamily separates in the signless-equality direction.
\end{corollary}
\begin{proof}
All cycle lengths are $6$ or $r+3$, so the twin-freeness argument
applies for $r\geq2$. The exceptional metric dimensions and all
remaining metric dimensions follow from
\cite[arXiv version, Theorems~8 and~14]{KSS}, with edge-length parameters
$(3,3,r)$: the exceptional multisets are exactly $(3,3,3)$ and $(3,3,5)$.
For completeness, the distance pairs to $(a,y_1)$ for $r=2$, in vertex
order $(u,v,a,b,c,d,y_1)$, are
\[
(1,1),(2,1),(0,2),(1,2),(2,2),(3,2),(2,0).
\]
For $r=4$, in order $(u,v,a,b,c,d,y_1,y_2,y_3)$, they are
\[
(1,1),(2,3),(0,2),(1,3),(2,2),(3,3),(2,0),(3,1),(3,2).
\]
Each list is injective. The same degree-three obstruction excludes
metric dimension one.

For all $r\geq3$, the nonsingular spectral argument still applies:
at $t=2e$ one has $P=e^{r-1}r\ne-3/2$; at the large quadratic conjugate
one has $|p_{r-2}(T)|>1$, since $r-2\geq1$.
The singular systems depend only on $r$ modulo six.
This proves the asserted table also for $r=3,4,5$.
At $r=2$, direct expansion of $\det(zI-L)$ and $\det(zI-Q)$ gives
\begin{align*}
 \chi_L(z)&=z(z-1)(z-3)(z^2-6z+7)^2,\\
 \chi_Q(z)&=(z-1)(z-3)(z^2-4z+1)(z^3-8z^2+17z-8).
\end{align*}
This length-two graph is already the second graph in
\cite[Table~2, p.~23]{AKS}, under the relabeling
$(u,v,a,b,c,d,y_1)\mapsto(1,4,2,3,6,5,7)$; its aggregate values
$(\dim,\mm_L,\mm_Q)=(2,2,1)$ are recorded there.
The quadratic in $\chi_L$ has roots $3\pm\sqrt2$, distinct from $0,1,3$.
The two nonlinear factors of $\chi_Q$ have discriminants $12$ and $316$,
respectively; the remainder of the cubic modulo the quadratic is $-4$,
and neither factor vanishes at $1$ or $3$.
Thus $\chi_Q$ is squarefree. Comparing the resulting maxima with the
metric dimensions proves all equality assertions.
\end{proof}

\begin{corollary}[A specialization of Paone--Paone (2026)]\label{cor:two}
For every $r\geq2$,
\[
 \operatorname{mult}_{Q(\Th_{3,3,r})}(2)
 =\begin{cases}1,&r\equiv1\pmod4,\\0,&\text{otherwise}.\end{cases}
\]
This is the $(3,3,r)$ instance of \cite[Corollary~8.1]{PP}.
\end{corollary}
\begin{proof}
Put $e=-1,t=0$ in \eqref{eq:short}--\eqref{eq:rec}.
Then $C=-A$, $B=-D_0$, $s=Y_1=-A+2D_0$, and the long-path
sequence is $A,s,-A,-s,\ldots$.
For $r\equiv0,1,2,3\pmod4$, the two remaining endpoint equations
reduce respectively to
\[
\begin{array}{c|cc}
r\bmod4&\multicolumn{2}{c}{\text{endpoint conditions}}\\\hline
0&D_0=A&-s=A\\
1&D_0=s&A=-D_0+2A\\
2&D_0=-A&s=3A\\
3&D_0=-s&-A=-D_0+2A
\end{array}
\]
Using $s=-A+2D_0$ leaves a one-dimensional solution space only
in residue one, and the zero space in the other three residues.
The antisymmetric subspace has eigenvalues $1,3$, not $2$.
\end{proof}

\section{A rooted-path consequence in both directions}
The restriction to $\Th_{3,3,r}$ precludes signless-only separation,
but a different construction gives both directions of Problem~30.
It combines the base examples in \cite[Theorems~10, 12 and Corollary~17]{AKS}
with known rooted-product results. The one-leaf metric criterion is
Jannesari's Theorem~2.2 \cite{J}; Feng and Wang
\cite[Proposition~2.10, equation (5)]{FW} extend it to longer paths
for diameter-two bases. The spectral block formula is
Cardoso--Martins--Robbiano--Rojo \cite[Theorem~2.1, equation (2.4)]{CMRR}.
We give the short specialization explicitly.

\begin{corollary}\label{cor:rooted}
For $t\geq1$, form $H_t$ by attaching one path of $t$ edges to every
vertex of $H$, identifying that vertex with one endpoint of its path.
If $H=K_{1,3,3}$ then
$(\dim H_t,\mm_L(H_t),\mm_Q(H_t))=(4,4,5)$.
If $H=K_2\vee(2K_2)$ then the triple is $(3,2,3)$.
Both constructions give infinite pairwise nonisomorphic families of
connected twin-free graphs. Thus Problem~30 has examples in both directions.
\end{corollary}
\begin{proof}
Write a vertex as $(v,j)$, where $v\in V(H)$ and $0\leq j\leq t$ is
its depth on the attached path. For a metric basis $B$ of $H$, take
the landmarks $\{(b,t):b\in B\}$. The distance to $(b,t)$ is $t-j$ if
$v=b$, and $j+t+d_H(v,b)$ otherwise. A vertex above a basis vertex
is identified by its unique coordinate at most $t$, then by that
coordinate's value. All remaining profiles are
$j+t+r_H(v\mid B)$. They are distinct provided no two outside
profiles differ by a nonzero constant vector.
For $K_{1,3,3}$ choose two vertices in each part of size three.
The outside profiles are $1111,2211,1122$; for $K_2\vee(2K_2)$ choose
one vertex from each of its three twin pairs. Its outside profiles
are $111,112,121$. Neither list contains a pair differing by a
nonzero constant vector. These bases have minimum sizes four and
three respectively, since a resolving set omits at most one vertex
from each twin class. Conversely, projection of any resolving set
of $H_t$ to its base vertices resolves the roots, because
$d((v,0),(b,j))=d_H(v,b)+j$. Thus $\dim H_t=\dim H$.

For $X=L,Q$, diagonalizing $X(H)$ decomposes $X(H_t)$ into the
blocks $X(P_{t+1})+\mu E_{00}$, one per base eigenvalue $\mu$,
with its multiplicity. Each block is irreducible tridiagonal and
has simple spectrum: an eigenvector's first coordinate determines
all subsequent coordinates by the row recurrence, and cannot be zero.
Its characteristic polynomial has the form $F(z)-\mu R(z)$,
where $R$ is the polynomial of the endpoint-deleted path block.
Consecutive endpoint minors are coprime: a common zero propagated
down the determinant recurrence would make the final constant
minor $1$ vanish. Hence $F$ and $R$ have no common zero, and blocks
for distinct $\mu$ have disjoint spectra. It follows that
$\mm_X(H_t)=\mm_X(H)$.
The base spectra are
\[
\begin{array}{c|cc}
H&\operatorname{Spec}L&\operatorname{Spec}Q\\\hline
K_{1,3,3}&0,4^{(4)},7^{(2)}&1,4^{(5)},9\\
K_2\vee(2K_2)&0,2,4^{(2)},6^{(2)}&2^{(3)},4^{(2)},8.
\end{array}
\]
These agree with the cited base formulas.
Both bases have minimum degree at least two, so the roots in $H_t$
have degree at least three and have distinct attached neighbors.
The nonroots have degree at most two and lie on pendant paths:
none lies on a triangle or quadrilateral, and distinct leaves have
distinct parents. These facts exclude adjacent and nonadjacent twins,
including a root paired with a nonroot. Connectivity is immediate,
and the orders $|V(H)|(t+1)$ distinguish the members of either family.
\end{proof}

\section{Coverage and accompanying checks}
For the modular-partition part of Problem~32, the following general
lower bound is already implicit in Belmonte--Fomin--Golovach--Ramanujan
\cite[Section~4, Lemma~7, p.~16]{BFG}.
For a proper module $M$, let $\rho(M)$ be the minimum number of
landmarks in $M$ separating its vertices using distances in $G[M]$
with every distance capped at two (including disconnected pairs).
\begin{corollary}\label{cor:module}
If $M_1,\ldots,M_s$ are pairwise disjoint proper modules in a connected
graph, then $\dim G\geq\sum_i\rho(M_i)$.
\end{corollary}
\begin{proof}
Every vertex outside a module has the same distance to all its members,
so only landmarks inside it can separate them. Connectivity makes
distances between its members equal to the capped induced distances.
Thus every resolving set has at least $\rho(M_i)$ members in each $M_i$;
disjointness permits summation.
\end{proof}
For a twin class of size $m$, this gives $\rho=m-1$.
The source's open problems are in \cite[Section~9, p.~27]{AKS}.
\begin{center}\small
\begin{tabular}{@{}p{27mm}p{118mm}@{}}
\toprule
Source item&Coverage\\\midrule
Problem 28&Partially answered: a twin-free mechanism, via short-path
symmetry and spectral collisions; the characterization of all connected
graphs is not covered.\\
Problem 29&Answered for $\Th_{3,3,r}$, every $r\geq2$:
Corollary~\ref{cor:complete}. Its congruence question has an affirmative
answer within this subfamily; general path lengths are not covered.\\
Problem 30&Answered: yes, by Theorem~\ref{thm:main}; no requirement for
both directions is imposed by the question. Corollary~\ref{cor:rooted}
also supplies both directions from prior rooted-product theorems.\\
Problem 31&Not covered: one graph per relevant order does not improve
the exponential lower bound of the source's Corollary~24.\\
Problem 32&Prior modular lower bound reproved as
Corollary~\ref{cor:module}; an equitable invariant-subspace example is
supplied below. General spectral-equality criteria are not covered.\\
Problem 33&Not covered: the theta graph has no pendant branches;
Corollary~\ref{cor:rooted} attaches only one path per root, rather
than the question's $p_i\geq2$ copies.\\
\cite[Table~2, row~2]{AKS}&Prior proof cited: the aggregate metric and
multiplicity values at $r=2$; the factorization above specifies its
repeated eigenvalues.\\
\cite[Theorem~25, Corollary~26]{AKS}&Prior proof cited: the twin-bearing
$\Th_{2,2,r}$ separation families; not a twin-free answer to Problem~30.\\
\cite[Theorems~8, 14]{KSS}&Reproved for the resolving-pair cases;
prior proof Knor--Sedlar--\v{S}krekovski (2022). The exceptional metric
dimensions are taken from their Theorem~8.\\
\cite[Theorem~4.4]{BFM}&Prior proof cited: the later alternative
metric-dimension proof of Benakli--Froitzheim--Martinez (2025), with
parameters $(2,2,r-1)$, after reordering.\\
\cite[Corollary~8.1]{PP}&Reproved specialization;
prior proof Paone--Paone (2026): multiplicity at $Q$-eigenvalue $2$.\\
\cite[Corollary~16]{BK}&Reproved; prior proof Bahmani--Kiani
(2015/2016): periodicity at $1,3$, not exclusion of other repeated roots.\\
\cite{J,FW,CMRR}&Prior results specialized and proved explicitly
in Corollary~\ref{cor:rooted}; the base spectra are from \cite{AKS}.\\
\cite[Section~4, Lemma~7]{BFG}&Reproved; prior proof
Belmonte--Fomin--Golovach--Ramanujan (2016): modular lower bound.\\
Signless-only requirement&Not posed for this subfamily; obstruction
proved by Corollary~\ref{cor:complete}.\\
\bottomrule
\end{tabular}
\end{center}
All Problems 28--33 in this table refer to \cite[Section~9, p.~27]{AKS}.
For Problem~32, the cells $\{a,c\},\{b,d\}$ and all remaining singleton
cells form an equitable partition for both matrices: it is the orbit
partition of the short-path swap. Its cell-constant subspace is exactly
the symmetric subspace above, while its orthogonal complement is the
two-dimensional antisymmetric subspace, with eigenvalues $1,3$.
Their collisions with symmetric eigenvalues produce the equalities
without twins. This decomposition is an example, not a general partition
lower-bound theorem.
Corollary~\ref{cor:module} supplies the separate modular lower bound.
The metric dimension also follows from \cite[arXiv version, Theorems~8 and~14]{KSS};
the explicit resolving pairs give direct certificates outside $r=3,5$.
The proof in Sections 2--3 establishes the full theorem and
Corollary~\ref{cor:all-lengths}. The accompanying \texttt{verify.py} checks
every integer $2\leq r\leq40$ by exact distance calculations and
polynomial greatest common divisors. In particular, for $1\leq k\leq6$:
\[
 \gcd(\chi_L,\chi_L')=(z-1)(z-3),\qquad
 \gcd(\chi_Q,\chi_Q')=1.
\]
For $r=3,5$ it exhausts all pairs and supplies a resolving triple;
otherwise it checks the displayed resolving pair. It also checks the
endpoint factorization, all twenty-four singular systems, the exceptional
length-two factors and Corollary~\ref{cor:two}. It also checks both bases
of Corollary~\ref{cor:rooted} for attached paths of one through four edges:
connectivity, twin-freeness, leaf resolving sets and exact spectral
multiplicities. These finite checks support,
but are not premises of, the
all-parameter proof.
The Lean declaration
\href{https://github.com/the-omega-institute/trureturing/blob/d9bb8a406783c5222849b798f3c2ab738f653d75/D5/S3/Combinatorics/SpectralGraph/TwinFreeThetaSeparation.lean}{\texttt{TwinFreeThetaSeparation.result}}
has as its conclusion the existence of an infinite pairwise nonisomorphic family of
connected twin-free graphs with exactly one equality. Its witnesses are
$\Th_{3,3,6k+4}$, $k\geq1$, and its conclusion does not specify repeated
eigenvalues or classify other lengths. The full multiplicity classification,
the distance formulas and Corollaries~\ref{cor:two}, \ref{cor:rooted} have the written proofs
above; the accompanying computations cover the stated finite range.

\begin{thebibliography}{10}
\bibitem{AKS} M.~R.~Alfuraidan, S.~Khan and L.~Sauras-Altuzarra,
\emph{Metric dimension and maximum eigenvalue multiplicity of graph
Laplacians}, arXiv:2610.26875v1 (2026).
\url{https://arxiv.org/abs/2610.26875v1}.
\bibitem{KSS} M.~Knor, J.~Sedlar and R.~\v{S}krekovski,
\emph{Remarks on the vertex and the edge metric dimension of 2-connected
graphs}, Mathematics \textbf{10} (2022), no.~14, 2411.
\url{https://doi.org/10.3390/math10142411}.
Theorem~14 is also in arXiv:2203.07335v1, p.~10.
\bibitem{PP} A.~Paone and M.~Paone,
\emph{Line-graph inertia of roses and generalized theta graphs},
version~1.0, 1~August 2026, Corollary~8.1, pp.~10--11.
\url{https://doi.org/10.5281/zenodo.21744051}.
\bibitem{BFM} N.~Benakli, N.~Froitzheim and D.~Martinez,
\emph{Metric dimension of generalized theta graphs},
arXiv:2510.12100v1, 14~October 2025, Theorem~4.4, p.~11.
\url{https://arxiv.org/abs/2510.12100v1}.
\bibitem{BK} A.~Bahmani and D.~Kiani,
\emph{Graph reduction techniques and the multiplicity of the Laplacian
eigenvalues}, arXiv:1512.08265v5 (2016), first posted 27~December 2015.
\url{https://arxiv.org/abs/1512.08265v5}.
\bibitem{LL} X.~Liu and P.~Lu, \emph{Laplacian spectral characterization
of dumbbell graphs and theta graphs}, arXiv:1407.5255v2 (2014),
Proposition~2.6, pp.~3--4.
\url{https://arxiv.org/abs/1407.5255v2}.
\bibitem{J} M.~Jannesari, \emph{On the metric dimension of corona
product of a graph with $K_1$}, Kyungpook Mathematical Journal
\textbf{63} (2023), 123--129, Theorem~2.2, pp.~125--126.
\url{https://doi.org/10.5666/KMJ.2022.63.1.123}.
\bibitem{FW} M.~Feng and K.~Wang, \emph{On the metric dimension and
fractional metric dimension for hierarchical product of graphs},
arXiv:1211.1432v1 (2012), Proposition~2.10, p.~6.
\url{https://arxiv.org/abs/1211.1432v1}.
\bibitem{CMRR} D.~M.~Cardoso, E.~A.~Martins, M.~Robbiano and O.~Rojo,
\emph{Copies of a rooted weighted graph attached to an arbitrary
weighted graph and applications}, Electronic Journal of Linear Algebra
\textbf{26} (2013), Theorem~2.1, equation (2.4).
\url{https://doi.org/10.13001/1081-3810.1681}.
\bibitem{BFG} R.~Belmonte, F.~V.~Fomin, P.~A.~Golovach and
M.~S.~Ramanujan, \emph{Metric dimension of bounded tree-length graphs},
arXiv:1602.02610v1 (2016), Section~4, Lemma~7, p.~16.
\url{https://arxiv.org/abs/1602.02610v1}.
\end{thebibliography}
\noindent\small\copyright\ 2026 Haobo Ma and Wenlin Zhang.
Licensed under Creative Commons Attribution 4.0 International (CC BY 4.0).
\end{document}
